\chapter{What Makes a Story?}
\section{A Video Your Phone Edited for You}
Pick up something you touch every day: your phone. Open its photo gallery.

Many galleries hide a feature called ``Memories'' or ``Highlights''. Quietly working in the background, it selects dozens of photographs and a few clips from last summer, adds music and transitions, and edits a ninety-second film. It sends you a notification titled ``Those Wonderful Days''.

You watch. A seaside sunset, smiling friends, ice cream, fireworks. As the music starts, you really feel something, perhaps even the beginnings of tears.

But you remember that summer. The two-hour queue for the water park, the fierce argument with a friend one afternoon, the misery of carsickness on the return coach. One radiant group photograph was taken just after the quarrel, when adults insisted, ``Come on, smile!'' None of that appears. The phone is not lying: every photograph is genuine; every filmed second happened. It merely selected, arranged, and added music.

The same days have two versions: your remembered hot, irritating summer with occasional bright moments, and the phone's relentlessly heartwarming summer. Both use identical material. They differ only in what was chosen, its order, what was concealed, and what was enlarged.

Here is our question: what constitutes a story? Why can identical events become completely different stories when selected, ordered, and told differently? More troublingly, if telling varies so much, how much ``truth'' remains inside a story?

\section{A Storyteller's Wooden Block}
Come into another scene.

It is evening in a teahouse in an old Sichuan town, or any traditional storytelling venue you imagine. A dozen bamboo chairs stand amid steam rising from lidded teacups. Tea, sweat, and roasted melon seeds mingle in the air. The low stage holds one table; behind it sits a storyteller with a folding fan and a wooden block, used to produce a startling crack.

He is telling the Three Kingdoms story. Cao Cao retreats defeated along Huarong Road, men and horses exhausted. Guan Yu blocks the pass on horseback, blade across his path. Somewhere along the way, the sound of cracking seeds has stopped. The storyteller snaps his fan shut, leans forward, and lowers his voice:

``Guan Yunchang grips his Green Dragon Crescent Blade and looks at Cao Cao. Kill him, or release him---''

The wooden block cracks like thunder.

``To learn what happens next, hear the next instalment!''

The room erupts. Someone shouts, ``Finish before you go!'' Others laugh, curse, and toss seed shells towards the stage. Smiling, the storyteller salutes with clasped hands, packs up, and leaves. Tomorrow at this hour, nearly everyone will return and pay for tea first.

Notice the strangest thing. The storyteller knows the ending, and the audience knows he knows. Both understand the arrangement. Yet refusing to finish the story does not drive customers away; it is the business's core technology. The teahouse sells not stories but the discomfort of a story half-told.

Why? If value lay only in ``what happened'', one sentence would suffice: Guan Yu remembered Cao Cao's past kindness and let him go. Ten Chinese characters, finished. Nobody would buy tea daily for those ten characters. They pay for something other than the events: suspension, waiting, and a pause precisely calculated to fall at ``next instalment''.

What is that ``something else''? It is our first clue.

\section{Between a Chronicle and a Story}
Lay out the conflict before reaching a conclusion.

Compare two pieces of writing. First, a diary:

``Up at seven. Breakfast at seven twenty: steamed buns. School at eight. Maths was third period. Lunch in the canteen. Ran eight hundred metres in afternoon PE. Homework in the evening. Bed at ten thirty.''

Everything is there: times, places, events, all true. Yet you would call it a list, not a story.

The second has one sentence: ``In the last hundred metres of the PE race, your desk-mate, a full half-lap behind, suddenly begins to sprint.''

One sentence, and your ears prick up. Why sprint? Did the runner catch up? What is the relationship between these two? You want more.

In \emph{Aspects of the Novel}, E. M. Forster offered a famous example. In essence, ``The king died, and then the queen died'' is a sequence of events, while ``The king died, and then the queen died of grief'' is a plot. The difference is the addition of ``of grief'', three characters in the Chinese phrasing. A list says what happened; plot says how things connect.

Over two thousand years earlier, Aristotle made the point in the \emph{Poetics}. Plot organises action, with a beginning, middle, and end. Events must follow not merely as ``this, then that'', but through probability or necessity. Storytelling's skill lies not in the pieces but in their connections.

Here comes the conflict. See clearly that it contains three questions, not one.

First, what is a story's material: events or their connecting lines? The chronicle camp says facts are material and connections are the mind's embroidery. The plot camp says unconnected events cannot enter the mind at all, like loose beads impossible to pick up together.

Second, are there multiple ways to connect them? The phone video has already answered yes. One summer can become heartwarming or humiliating. Whoever makes those connections holds formidable power: deciding whether you laugh or cry and whom you regard as good, without changing a single event.

Third, if telling varies so greatly, can stories still be ``true''? The phone forged no photograph; it merely chose. Yet would you call its warm little film ``the truth about that summer''?

We will take these questions one by one. First, decide which feels truer: the messy summer in your memory or the warm summer edited by your phone. Remember your answer; we will return to it.

\section[Who Wants to Know the Ending?]{Those Who Want the Ending and Those Who Cover Your Ears}
The teahouse contains at least two kinds of listeners. Both deserve a full hearing.

\textbf{The first cannot wait to know the ending.}

You know such people. They turn to a novel's last page, search for a television finale, or interrupt a story with ``Then what? Just tell me the result.'' Before criticising them, hear their complete case.

First, suspense is a tax, not a gift. Anxiety over whether someone dies runs like a background program, consuming memory and preventing appreciation of the present passage. Know the ending and the mind is free to notice craft: the sentence's construction, the planted clue. Second readings of \emph{Dream of the Red Chamber} are often more attentive because readers are no longer rushing.

Second, this may be more than taste. In 2011, two American psychologists published experiments in the influential journal \emph{Psychological Science}. Some readers received short stories with their endings revealed at the beginning. Contrary to intuition, their average enjoyment did not decline; in several groups it rose slightly. Later research tested and qualified the result, so it is no iron law. But \textbf{the supposedly obvious rule that spoilers ruin stories is less solid than it seems}. This is a classic easily misjudged case.

Third, life is uncertain enough. Why require uncertainty in entertainment? Some people watch films to relax, not suffer.

\textbf{The second kind fiercely rejects spoilers.}

Their case is equally strong. A first encounter is nonrenewable. You have only one first experience of a twist; once burned, it is gone, whereas you can always choose to burn it later. Moreover, the sequence of expectations is part of the work. Striking the block at ``kill or release'' is no arbitrary choice but the performance's centre of gravity. A spoiler does more than reveal an ending: it \textbf{makes an unauthorised decision for someone else}, turning their first experience into your second-hand material.

This is not impatience versus patience. It is two answers about stories' value. One places value in information, the sooner acquired the better. The other places it in experience, which has its own route and cannot skip stops.

Pull back to see suspense's weight in another civilisation.

You probably know the frame of \emph{The Thousand and One Nights}. Betrayed, King Shahryar comes to hate women, marrying a girl each day and killing her at dawn. The vizier's daughter Scheherazade volunteers to enter the palace. Her only weapon is storytelling. Each night she tells a captivating tale and stops at its most urgent moment when dawn arrives. Wanting the next part, the king lets her live another day. The following night she finishes it, begins another, and again stops at the crucial moment. So it continues for a thousand and one nights.

Weigh that frame. At its extreme, ``hear the next instalment'' is not salesmanship but survival. Every pause must be exact: too early and curiosity is insufficient, too late and there is no chance to stop. The storyteller's block and the vizier's daughter's silence are the same instrument: a precise calculation of how long curiosity can wait.

Leave room for a third voice, unpleasant but enduring: \textbf{``Stories are invented and false. Listening wastes time and may even mean being deceived.''}

Give it its full case. Reality is complex enough: famine, war, equations, experimental data all deserve effort. Is it not absurd to shed real tears over nonexistent people or lose sleep over invented conflict? Worse, stories naturally persuade. Advertising, propaganda, and rumours wear their clothes because stories bypass rational security checks. Someone raised only on stories, never data, is easy to deceive.

Each reason has force, making it important to locate the error: confusing fiction with lying. A lie claims something false is true. A story begins with a contract: ``Once upon a time'', ``Let me tell you a story''. Speaker and listener understand that imagination, not reporting, follows. Scheherazade's king is not deceived, nor is anyone in the teahouse. Stories are among the few human inventions that openly announce invention and still attract willing listeners. Why they deserve listening to despite being invented is precisely what remains to be answered.

\section{Thinking Bridge: Two Maps of One Story}
After all this, take away a tool. It is a foundational move in narratology, the study of how stories are told: \textbf{draw two maps of the same story.}

\textbf{Map one: the order in which events happen.}

This is the story world's chronological timeline. For Mulan taking her father's place in the army: the khan calls up troops $\rightarrow$ Mulan decides to go for her father $\rightarrow$ buys horse and saddle $\rightarrow$ leaves home $\rightarrow$ fights for ten years $\rightarrow$ returns victorious $\rightarrow$ declines office $\rightarrow$ goes home $\rightarrow$ resumes women's clothing.

\textbf{Map two: the order in which narration presents events.}

This is the order in which the audience receives information: what the author shows first and later, skips, or tells twice.

The gap between the maps contains all the possibilities of telling. Use four steps:

First, \textbf{take it apart}. Split the story into its smallest units, events. Mark four things in each: characters, who is present; desire, who wants what; obstacle, what blocks them; and change, what differs afterwards. Almost every story across times and cultures contains these. Without desire, characters are scenery. Without obstacles, desire is immediately satisfied and the story ends in three sentences. Without change, listeners ask why they bothered when the ending is the starting point.

Second, \textbf{arrange map one}. Put the pieces in the story world's temporal order.

Third, \textbf{mark map two}. Identify the narrator's actual treatment, principally five devices:

\begin{itemize}
\item \textbf{Suspense}: delay the answer. ``Kill or release? Hear the next instalment.'' Suspense hides answers.
\item \textbf{Foreshadowing}: plant a component early. Chekhov repeatedly argued in letters and conversation that a gun hanging on the wall in the first act should fire later; otherwise, do not hang it. Foreshadowing works opposite to suspense: it hides not the answer but the fact that something is a question. Nobody initially knows the gun matters.
\item \textbf{Flashback}: show later events first and supply earlier ones afterwards. The \emph{Odyssey} opens with Odysseus already trapped on Calypso's island for seven years. His preceding ten years of wandering are recounted later at someone else's court. The ancient Roman poet Horace noticed Homer's technique over two thousand years ago, praising his plunge into the middle of things. Films opening with an ending and then a black screen reading ``Three years earlier'' use the same move.
\item \textbf{Omission}: cut an entire stretch. The \emph{Odyssey} does not describe Odysseus growing up; the \emph{Ballad of Mulan} omits ten years of warfare. We will examine that cut next.
\item \textbf{Repetition}: tell something repeatedly. In Homer, dawn is always rosy-fingered and the sea wine-dark. In oral performance, formulas anchor memory and improvisation. Chinese storytelling's ``Two flowers bloom; let us follow each branch'' similarly institutionalises repetition and switching.
\end{itemize}
Fourth, \textbf{ask why}. Overlay the maps. Every mismatch raises a question: why begin in the middle, cut this passage, or plant that detail so early? Telling is not decoration; it is part of the work's content.

Practise with a tiny story, Aesop's tortoise and hare. Characters: tortoise and hare. Desire: winning. Obstacles: the hare's pride and the tortoise's short legs. Change? Think carefully. Neither animal really changes; the hare may remain proud next time, and the tortoise still has short legs. \textbf{The listener is what truly changes.} This is a fable's secret: it relocates change from people inside the story to people outside it.

Try both maps for any novel, film, game, or even news item. Vague feelings that a story is well or badly told often become specific, explainable judgements once mapped.

\section{The Twelve Years Mulan Cuts Out}
Read a short primary text and see how forcefully it uses the two maps.

The \emph{Ballad of Mulan}, a Northern Dynasties folk song preserved in the Song-compiled \emph{Collection of Music Bureau Poetry}, tells of Mulan replacing her father in the army. You know the story; today we examine its telling.

It begins:

\begin{quote}
Sigh upon sigh: Mulan weaves beside the door. No loom or shuttle sounds; only the young woman's sighing.
\end{quote}
Four lines supply the four components: Mulan, the conscription order and ageing father, and the conflict between desire and obstacle contained in a sigh.

She decides to leave, and prepares:

\begin{quote}
At the eastern market, buy a fine horse; at the western market, a saddle and pad; at the southern market, a bridle; at the northern market, a long whip.
\end{quote}
Pause. On map one, must buying a horse and equipment require four markets? One would suffice. Yet narration expands the act into four sentences and four directions. This is \textbf{repetition and elaboration}. It serves feeling, not efficiency. Following her four trips makes you feel the bustle, solemnity, and gravity of what is happening.

Then come the poem's boldest two lines:

\begin{quote}
Generals die in a hundred battles; warriors return after ten years.
\end{quote}
Ten years. A hundred battles. Mortal danger. Fourteen Chinese characters, finished.

On map one, warfare occupies at least ninety percent of those actual ten years. On map two it occupies two lines. This is not laziness. Look at her return home:

\begin{quote}
Open my eastern chamber's door, sit upon my western chamber's bed. Take off my wartime robe, put on my old clothes.
\end{quote}
Ten minutes at home receive elaborate treatment again. Open the east door, sit on the west bed, slowly change clothes, sentence by sentence.

Overlay the maps and the answer appears: the poem's allocation of time declares its theme. It does not care principally about war, which it dismisses with a stroke. It cares about a woman entering a man's world, coming back, and becoming herself again. Hence the attention to dressing before departure and changing after return. A final comparison closes it:

\begin{quote}
The male rabbit's feet fidget; the female rabbit's eyes are blurred. When two rabbits run beside each other, how can you tell whether I am male or female?
\end{quote}
\textbf{Omission is not cutting corners; it tells you what the author's story is about.} What narration removes or enlarges can reveal its position more clearly than what it says. Remember this when we turn to today's news.

\section{Why Is Cinderella Everywhere?}
Now face a fact requiring explanation.

You may think Cinderella comes from Europe. Perrault's 1697 French version supplied glass slippers; the Grimms' 1812 German version made the stepsisters pay a bloody price. Yet centuries earlier, Tang writer Duan Chengshi's \emph{Miscellaneous Morsels from Youyang} recorded Ye Xian. After her father's death, her stepmother mistreats her and kills and eats her magical fish. Ye Xian hides its bones, which grant her wishes. Wearing golden shoes, she secretly attends a festival, flees when recognised, and loses a shoe. Eventually it reaches a king, who has women throughout the kingdom try it on and finds its owner.

Stepmother, abused girl, magical helper, shoe, identification by shoe: almost the same framework as European Cinderella. This is only one family of similar stories. The oppressed rise, heroes leave for adventures and return, the youngest of three brothers succeeds, the greedy are punished. Such tales span Europe, the Middle East, India, China, and Japan. Why?

Separate four matters first. \textbf{What happened}: the stories exist and their frameworks resemble one another, verifiable and undisputed. \textbf{How people then understood them}: Ye Xian's magical fish and recompense probably suggested to Tang listeners that goodness receives rewards. \textbf{How we evaluate them}: that is for you. What follows compares the explanatory layer, \textbf{why this happened}. At least three accounts offer different reasons and limitations.

\textbf{Explanation one: stories travel.}

Similarity comes from shared ancestry. Stories accompany traders, armies, migrants, and marriages along roads, changing into local clothes at each stop. A fish becomes a fairy godmother and golden shoes become glass, but the framework remains. The argument is strong: details align too closely, especially the specific device of finding someone by trying on a shoe, to make independent invention likely. Eurasian trade routes certainly existed; stories can follow any road goods travel.

Its limits are clear too. First, it cannot explain similar stories in places with little possible contact, such as flood myths on several disconnected continents. Second, it answers how stories travelled, not the earlier question of why these particular types survived among thousands capable of travelling. Transmission requires transmissibility, which diffusion alone cannot explain.

\textbf{Explanation two: minds share story templates.}

This account looks inside the mind. In \emph{The Hero with a Thousand Faces}, twentieth-century American scholar Joseph Campbell proposed that peoples' heroic legends share a framework: the hero leaves ordinary life, enters danger, undergoes trials, and returns with a reward, the ``hero's journey''. Cinderella and Ye Xian resemble one another not because one copied the other but because human minds have this shape. Cultures are different factories independently producing similar forms.

Grand in scope, the explanation has considerable limits. An abstract template can fit anything. School becomes departure from ordinary life, an examination a trial, summer holiday a return. An account that explains everything often explains nothing, because no possible case escapes it. Campbell has also long been criticised for trimming important differences between stories to fit a common frame.

\textbf{Explanation three: stories are simulators for the brain.}

This digs deeper. Human brains naturally seek causes. Our ancestors survived by connecting encounters into chains of who did what, wanted what, and obtained what result: who stole whose prey, who helped whom, where danger appeared. Such information naturally takes the shape of characters, desire, obstacle, and change. Stories shaped this way are easiest to understand, remember, and transmit. People abused or dismissed especially need stories of reversal, free rehearsals in which those similarly placed can first win imaginatively. Cinderella is everywhere because belittled people are everywhere.

This explains both transmissibility and form, apparently a complete account. Its problem is how easily it can be told. Almost anything that survives can receive a polished retrospective story about its usefulness to ancestors, difficult to test. In this chapter's terms, beware of making an explanation too smooth a story.

The three accounts are distinct without wholly conflicting. Diffusion asks how stories travelled, templates why they have this shape, simulation why minds accept them. The methodological reminder is simple: \textbf{when an explanation explains everything, check whether it discarded something in advance.}

\section[Further Challenge: A Hundred Fairy Tales]{Further Challenge: Taking a Hundred Fairy Tales Apart}
Here is our weightiest theory, far more radical than two maps.

In 1928, Russian scholar Vladimir Propp published the book later translated as \emph{Morphology of the Folktale}. His method sounded laborious: dismantle a hundred Russian wonder tales and compare events one by one to count kinds of action. The result astonished. Whether the protagonist was prince or shepherd, the antagonist witch or dragon, actions fell into only thirty-one ``functions'': departure, prohibition, violation, villainous reconnaissance, acquisition of a magical object or helper, confrontation, punishment, marriage or enthronement, and so on. Their order remained broadly fixed across the stories. Characters, seemingly endlessly varied, fell into seven roles by function: hero, villain, donor, helper, sought person, dispatcher, and false hero.

Weigh the discovery. Characters can change while functions stay put. Replace Ivan with a prince and a dragon with a witch; the machine and components remain, only the paint changes. Propp effectively declared: \textbf{who matters less than which actions happen in which order.}

This immediately explains familiar experiences. Why can you predict some films' endings after ten minutes? An internal function chart runs automatically: ``The false hero has arrived; exposure comes next.'' Why does a fifth sequel feel like the same medicine in new packaging? The function sequence remains, only its shell changes. Hollywood screenplay manuals and game quest design have Propp in their bones, whether they know his name or not.

Its sharpest use answers our opening question: \textbf{how can one story become comedy, tragedy, detective fiction, or fable?} Narratology distinguishes what is told, the events or functional layer, from how it is told, the narrative layer. Functions can remain identical while altered telling changes genre. Test it with three events:

\begin{quote}
Late at night, a daughter comes home to find her mother in the unlit sitting room. A sheet of paper lies on the table: the report card she secretly threw away that day.
\end{quote}
Keep the material fixed, adding nothing, and change only the telling:

\textbf{Detective story}: begin with the report card's disappearance. The daughter put it in the bin downstairs. How did it reach the table? Who retrieved it? Where has her mother been? The reader follows her backwards through doubts. Detective fiction reverses the causal chain, showing effects first and excavating causes bit by bit.

\textbf{Tragedy}: begin earlier that day. At the parents' meeting the mother forces a smile; on the way home she stands beside the bin for a long time. The reader knows everything, the daughter nothing. You watch her enter, switch on the light, and open her mouth to invent a lie, unable to intervene. The central device is that \textbf{the audience knows more than the character} and can only watch the train approach the collision.

\textbf{Comedy}: begin with the daughter and younger brother blaming each other. ``He picked it up!'' ``She threw it away!'' Misunderstandings multiply; explanations worsen matters. Finally the mother takes a yellowed sheet from a drawer: her own report card, discarded thirty years earlier. Everyone dissolves into laughter. Comedy's device is mismatch and release: everyone performs the wrong script, but the cost is low enough for a safe landing.

\textbf{Fable}: strip detail to the bones. ``A child threw bad news into the river. That night, the bad news walked home.'' Add a moral: hidden things grow legs and return. Fables push omission to the limit, removing the specifics of this time and retaining the shape of every time.

Four genres, the same events. The differences belong to narration: where to begin, what to hide or show, whose position the reader occupies, how far to open information's gates. \textbf{Genre lies in telling, not material.} Thus what a story recounts and how well it is told are separate questions. Ordinary material can become gripping, and extraordinary events can be told as dull as chewed wax.

Finally, give this theory the same scrutiny as others. Propp's chart came from a hundred Russian tales; beyond that forest it does not always work. Fit its thirty-one functions to \emph{Dream of the Red Chamber} and it repeatedly jams. More importantly, it misses something central. Two stories with identical function sequences can make your scalp prickle or make you yawn. The difference lies in tone, viewpoint, whose eyes see and whose voice speaks, all outside the chart. Later narratology pursues precisely this: not merely what action happens but who sees, who speaks, and where the reader is placed. Propp supplies a skeleton; other tools must supply flesh.

\section[Trending Stories and Life Edits]{Today: Trending Stories, Short Videos, and Life Edits}
These ancient questions now operate inside your pocket.

Start with short videos. You have seen ``a whole film in three minutes'', a rapid voice compressing a hundred-minute story into one map: ``Watch closely; this man is called Xiaoshuai.'' Efficient, certainly: all the functions in three minutes. But calculate the exchange. You receive map one and lose map two, where the work's pacing, heartbeat, and precise strike of the wooden block live. A recap never gives you that suspended instant of ``kill or release''. Recaps are not guilty; simply know what you bought: information, not experience. Both are useful, but do not pass one off as the other.

Now consider web-fiction and wish-fulfilment drama production lines. Engagement cancelled $\rightarrow$ unexpected opportunity $\rightarrow$ public humiliation of an opponent $\rightarrow$ a stronger patron appears $\rightarrow$ another public reversal. Recognise it? Propp's functions on a round-the-clock assembly line. Do not rush to sneer. Homeric formulas and the storyteller's ``next instalment'' were production lines too. Repetition has always belonged to the story industry. The danger is not formulas but that \textbf{when every story on the shelf shares one skeleton, only one voice remains}: revenge gratified, fortunes reversed instantly. Nobody tells the slow, silent days of overlooked people whose fortunes never reverse. Templates determine not only entertainment preferences but which circumstances society can see.

Most important is news. Three outlets can narrate one event as three genres. A makes tragedy, beginning with childhood and adding solemn music, leaving you only a sigh. B makes detective fiction: surveillance, doubts, timelines, ``Is it really that simple?'', leaving you waiting for a twist. C's comment section makes a fable, removing circumstances and retaining ``I always said people must never\ldots'' Facts can be identical while stories differ completely. This is self-defence, not literary appreciation. Before taking sides on an enraging trending story, ask where the telling begins, whom it makes protagonist or background, and what it omits. Remember Mulan's twelve missing years? \textbf{Omission reveals a position most honestly.}

Look at yourself too. Social posts, updates, and year-end reflections are edited works. An algorithm edits your gallery's memories; you edit your account. Which photograph, which order, which caption, which days never appear? We mock the phone's false edit while editing even more fiercely ourselves. That is not hypocrisy but the human default: we live by stories, not chronology.

Finally, a new variable: machines that write stories. Today's chatbots have read more tales than any person, know all thirty-one functions, and can draw both maps. What they lack is a lived situation. They do not know why you want this story tonight in this mood. Unlike Scheherazade, they do not die if they tell it badly. Stronger tools make who tells a story, to whom, and why now increasingly our responsibility.

\section{For You: Does Your Life Have a Plot?}
Return to the phone video.

At the start you judged whether the messy or heartwarming summer was truer. Expand the question to your life: is it a list of events or a plotted story?

If you told a story of the past year, where would you begin? Who is the protagonist? What desire and obstacle, and what changes? What would you cut?

The difficulty is that editing omits, and the omitted part is you too. The carsick, quarrelling, mistake-making self and the smiling photograph self: which is more real? Psychological research on self-understanding suggests that people know themselves by organising memories into stories. ``I am someone who gets up after falling'' is not a factual statement but a plot template used to select and arrange memories, then lived by. In other words, \textbf{you do not merely have stories; you are someone who has narrated yourself into your present form.}

Here is the final question, with no standard answer:

When fidelity to facts conflicts with telling a good story, whether describing your growth, a friendship's end, or yourself to new classmates, which should you choose? Is a chronicle preserving every embarrassment closer to the real you, or an honestly edited story that leaves embarrassment out?

Give your answer and reasons. First review your tools: two maps, four components, five devices, and Mulan's twelve missing years. No standard answer exists, but yours will have a shape: the story about yourself that you are telling now.
